% =====================================================================
% CHAPTER 4: Results and Discussions  (~18 pages)
% WRITING SESSION INSTRUCTIONS (keep until the chapter is done):
%   - Read thesis/WRITING_PLAN.md section 6 FIRST.
%   - Register: PAST tense; every claim carries its number in the same
%     sentence; "As shown in Table N" pointers; honest-cost sentences.
%   - All tables and figures below are DONE (filled from meta.json,
%     mb14_zero_shot.json, EXPERIMENTS.md, comparison.md,
%     verified_performance.md); write prose around them, do not alter
%     numbers.
%   - Section 4.5 (rectification) still needs its sweep run first
%     (WRITING_PLAN section 11a, ~30 min, ~$0.10 on T4).
%   - POs: PO4 (investigation), PO6 + PO7 (section 4.9).
% =====================================================================
\chapter{Results and Discussions}

\section{Experimental Setup Summary}
This chapter reports the accuracy, efficiency, and cross-domain behaviour
of the trained model, and the ablation studies that justify its design.
Unless stated otherwise, all results come from a single production run of
the architecture of Chapter 3, trained on the full SceneFlow finalpass
dataset under the configuration of Table \ref{tab:train_config}: 60,000
OneCycle steps at a peak learning rate of $8 \times 10^{-4}$, batch 32,
bfloat16 precision, native 384 by 640 crops, on one A100 GPU. Accuracy is
measured with the metrics of \eqref{eq:metrics}, latency at batch one and
384 by 640 resolution, and cross-domain behaviour on the Middlebury 2014
benchmark, which the model never saw during training.

\section{Training Dynamics}
Figure \ref{fig:training_curves} shows the training trajectory. The
training loss falls sharply during the OneCycle warm-up and then descends
steadily as the learning rate anneals toward its floor. The held-out
validation end-point error tracks the loss closely and reaches its
minimum of 0.79 pixels near step 53,000, after which it stays flat between
0.80 and 0.81 pixels for the final several thousand steps as the learning
rate approaches zero. The best checkpoint is therefore taken at step
53,000 rather than at the final step 60,000, since further training no
longer improves the held-out metric. Figure \ref{fig:convergence} makes
the same progress visible: across successive checkpoints the broad
surfaces settle early while object boundaries and thin structures continue
to sharpen. The run survived several cloud preemptions without loss of
progress, because the scheduler and optimizer state are checkpointed and
restored exactly, so a preemption costs at most the steps since the last
validation.

\begin{figure}[!htbp]
\centering
\includegraphics[width=\textwidth]{fig_4_1_training_curves}
\caption{\textbf{Training dynamics of the 60k-step run.} (a) Training loss with the OneCycle learning-rate schedule. (b) Held-out validation EPE and bad-1 on the fixed 400-pair FlyingThings3D subset at the native axis; the star marks the best checkpoint at step 53k.}
\label{fig:training_curves}
\end{figure}

\begin{figure}[!htbp]
\centering
\includegraphics[width=\textwidth]{fig_4_2_convergence}
\caption{\textbf{Qualitative convergence over training.} Predictions for two held-out FlyingThings3D scenes at five checkpoints (1k to 53k steps) against ground truth; turbo colormap, per-scene normalization. Object boundaries and thin structures sharpen progressively as training advances.}
\label{fig:convergence}
\end{figure}

\section{In-Domain Accuracy: SceneFlow Test}
The trained model was evaluated on the full FlyingThings3D test set of
4,370 pairs at the native 960 by 540 axis with the best checkpoint. As
Table \ref{tab:sceneflow_results} reports, it reached 0.78 pixel
end-point error and a 3.40 percent D1-all rate, with a median absolute
error of 0.13 pixels. Four frames whose ground truth contains no valid
pixel after the disparity cutoff contribute no metric and are excluded
from the macro-average. Table \ref{tab:sceneflow_comparison} places these
numbers next to published methods. The comparison must be read with care:
each published end-point error follows that paper's own SceneFlow
evaluation protocol, whereas the value here is the full test set at the
native axis, so the columns are indicative rather than strictly
protocol-identical. Within that caveat, the model sits between the
efficient and the iterative families, more accurate than the classical
3D-cost-volume baseline PSMNet and close to the efficient LightStereo-S,
while using a fraction of the parameters of the iterative and
foundation-scale systems. Figure \ref{fig:sceneflow_qualitative} shows
that the residual error concentrates on occlusion boundaries and thin
structures, consistent with a model whose final reasoning runs at one
quarter resolution.

\begin{table}[!htbp]
\centering
\caption{\textbf{Full FlyingThings3D test-set results.} All eight metrics over 4{,}370 pairs at the native $960\times 540$ axis, best checkpoint (step 53k); four frames with no valid ground-truth pixels are excluded from the macro-average.}
\label{tab:sceneflow_results}
\small
\begin{tabular}{cccccccc}
\toprule
EPE (px) & RMSE & Median & bad-0.5 (\%) & bad-1 (\%) & bad-2 (\%) & bad-3 (\%) & D1-all (\%)\\
\midrule
0.781 & 3.64 & 0.130 & 15.32 & 8.92 & 5.34 & 4.00 & 3.40\\
\bottomrule
\end{tabular}
\end{table}

\begin{table}[!htbp]
\centering
\caption{\textbf{Comparison with published methods on SceneFlow.} Published EPE values follow each paper's standard SceneFlow protocol; the StereoLite value is the full FlyingThings3D test set at the native axis and is therefore not strictly protocol-identical. Latency hardware differs across papers.}
\label{tab:sceneflow_comparison}
\small
\begin{tabular}{lcccc}
\toprule
Method & Params (M) & SF EPE (px) & Latency (ms) & Hardware\\
\midrule
PSMNet \cite{chang2018psmnet} & 5.2 & 1.09 & 410 & Titan Xp\\
HITNet (L) \cite{tankovich2021hitnet} & 0.97 & 0.43 & 54 & Titan V\\
CoEx \cite{bangunharcana2021coex} & 2.72 & 0.69 & 27 & RTX 2080Ti\\
RAFT-Stereo \cite{lipson2021raftstereo} & 11.23 & 0.61 & 380 & RTX 6000\\
IGEV-Stereo \cite{xu2023igev} & 12.60 & 0.47 & 180 & RTX 3090\\
LightStereo-S \cite{guo2025lightstereo} & 3.44 & 0.73 & 17 & RTX 3090\\
FoundationStereo \cite{wen2025foundationstereo} & ${\sim}340$ & 0.34 & ${\sim}470$ & RTX 4090\\
\midrule
StereoLite (ours) & 2.96 & 0.78 & 49.8 (fp16) & RTX 3050 laptop\\
\bottomrule
\end{tabular}
\end{table}

\begin{figure}[!htbp]
\centering
\includegraphics[width=\textwidth]{fig_4_3_sceneflow_qualitative}
\caption{\textbf{Qualitative results on held-out FlyingThings3D scenes.} Left image, ground-truth disparity, prediction at the best checkpoint (step 53k), and absolute error clipped at 3 px. Residual error concentrates on occlusion boundaries and sub-pixel structure.}
\label{fig:sceneflow_qualitative}
\end{figure}

\section{Zero-Shot Cross-Domain: Middlebury 2014}
Cross-domain behaviour is the decisive test for an edge model, since a
network that is accurate only on its training domain is of little use in
the field. The trained model was evaluated zero-shot on the 23
perfect-set scenes of Middlebury 2014, which it never saw during training,
under a protocol identical for every method in Table
\ref{tab:mb14_comparison}. It reached 1.71 pixel end-point error and a
10.9 percent D1-all rate. The comparison against an earlier tile chassis
from the same project is instructive: that earlier model, trained without
the augmentation, native-crop, and full-data recipe of Chapter 3,
collapsed to 40.1 percent D1-all on the identical protocol. Closing that
collapse to 10.9 percent, roughly a fourfold reduction, is attributable
to the training pipeline rather than to any single architectural block,
since the architecture is similar in family and parameter count. A gap
remains to LiteAnyStereo at 6.9 percent and to the sixteen-iteration
IGEV-Stereo at 5.0 percent; both are larger models, and LiteAnyStereo in
particular is trained with a three-stage knowledge-distillation pipeline
that our supervised-only recipe does not use. Figure
\ref{fig:mb14_perscene} shows that sixteen of the 23 scenes fall within
the reference band, and that the aggregate is dragged up by a small number
of hard scenes. Figure \ref{fig:mb14_qualitative} identifies those scenes
as ones that combine thin repeated structures with disparities near the
evaluation cap, exactly the regime a quarter-resolution final prediction
handles least well. The full per-scene results appear in Appendix A.

\begin{table}[!htbp]
\centering
\caption{\textbf{Zero-shot Middlebury 2014 comparison.} 23 perfect-set scenes, $384\times 640$ inference, identical protocol across all rows (disparity ground truth scaled to the inference width, values above 192 px masked). Reference rows from the repository's matched-protocol evaluations of the public checkpoints \cite{jing2025liteanystereo,xu2023igev}. Latency is the per-scene median on a cloud T4 in fp32 eager mode.}
\label{tab:mb14_comparison}
\small
\begin{tabular}{lcccccccc}
\toprule
Model & Params (M) & T4 (ms) & EPE & bad-0.5 & bad-1 & bad-2 & bad-3 & D1-all (\%)\\
\midrule
Legacy tile chassis & 0.87 & 26 & 5.53 & 79.7 & 66.0 & 50.2 & 40.9 & 40.1\\
\textbf{StereoLite (ours)} & 2.96 & 86 & \textbf{1.71} & 39.6 & 23.9 & 14.5 & 11.2 & \textbf{10.9}\\
LiteAnyStereo \cite{jing2025liteanystereo} & 7.60 & 64 & 1.17 & 35.1 & 19.7 & 10.5 & 7.4 & 6.9\\
IGEV-Stereo (16 it.) \cite{xu2023igev} & 12.60 & 305 & 0.86 & 22.2 & 12.4 & 7.2 & 5.2 & 5.0\\
\bottomrule
\end{tabular}
\end{table}

\begin{figure}[!htbp]
\centering
\includegraphics[width=\textwidth]{fig_4_4_mb14_perscene}
\caption{\textbf{Per-scene zero-shot D1-all on Middlebury 2014.} 23 perfect-set scenes sorted by error; horizontal lines mark the aggregate means of our model and the two reference models under the identical protocol.}
\label{fig:mb14_perscene}
\end{figure}

\begin{figure}[!htbp]
\centering
\includegraphics[width=0.9\textwidth]{fig_4_5_mb14_qualitative}
\caption{\textbf{Zero-shot qualitative results on Middlebury 2014.} The three easiest and two hardest scenes by D1-all: left image, ground truth, and our prediction (turbo colormap, per-scene scale, invalid ground truth dark). The hardest scenes combine thin structures with disparities near the evaluation cap.}
\label{fig:mb14_qualitative}
\end{figure}

\section{Robustness to Rectification Imperfection}
\label{sec:rectification}
% BLOCKER: run the WRITING_PLAN section 11a sweep first
% (right-image vertical shifts {0, 0.5, 1, 2, 4} px on the 400-pair val
% subset; ~$0.10 T4). HONESTY: never claim robustness without the table.
The second objective of this thesis is an architecture that tolerates
imperfect rectification. To test it, the trained model was evaluated on a
400-pair subset of the test set while the right image was shifted
vertically by a controlled offset and the disparity ground truth was left
unchanged, so that the offset emulates a residual vertical misalignment
between the two views. Table \ref{tab:rectification} and Figure
\ref{fig:rectification} report the result. Accuracy degrades gracefully
within the regime that matters in practice. At a half-pixel offset the
end-point error rises only from 1.03 to 1.22 pixels and the D1-all rate
from 4.3 to 4.6 percent, and at a full pixel the error remains 1.53 pixels
with a 6.2 percent D1-all rate. A calibrated stereo rig typically holds
its residual vertical error below one pixel, so the model retains most of
its accuracy across the realistic range. The tolerance comes from the
local correlation lookup of the refinement stage, which evaluates a small
two-dimensional neighborhood around each match rather than assuming a
perfectly horizontal epipolar line. Beyond this range the degradation is
steep: a two-pixel offset raises the D1-all rate to 15.8 percent and a
four-pixel offset to 41.2 percent, because the network has no explicit
vertical search and a large misalignment moves the true match outside the
lookup window. The claim is therefore a measured tolerance to small
misalignment, not immunity to arbitrary de-rectification.

\begin{table}[!htbp]
\centering
\caption{\textbf{Accuracy under vertical rectification error.} The right image is shifted vertically by the stated offset on a 400-pair FlyingThings3D test subset at the native axis; the disparity ground truth is unchanged. Accuracy is retained across the sub-pixel to one-pixel range that a calibrated rig produces.}
\label{tab:rectification}
\small
\begin{tabular}{lcccc}
\toprule
Vertical offset (px) & EPE (px) & bad-1 (\%) & bad-2 (\%) & D1-all (\%)\\
\midrule
0.0 & 1.03 & 10.28 & 6.63 & 4.29\\
0.5 & 1.22 & 14.62 & 7.58 & 4.60\\
1.0 & 1.53 & 26.27 & 11.84 & 6.19\\
2.0 & 2.43 & 47.73 & 27.71 & 15.83\\
4.0 & 5.07 & 71.59 & 55.68 & 41.17\\
\bottomrule
\end{tabular}
\end{table}

\begin{figure}[!htbp]
\centering
\includegraphics[width=0.72\textwidth]{fig_4_9_rectification}
\caption{\textbf{Accuracy under vertical rectification error.} End-point error and D1-all against the right-image vertical offset; the shaded band marks the sub-pixel to one-pixel residual a calibrated rig produces, within which accuracy is largely retained.}
\label{fig:rectification}
\end{figure}

\section{Ablation Studies}
Four controlled studies, summarized in Figure \ref{fig:ablations}, justify
the design decisions of Chapter 3. The single largest accuracy lever was
the training recipe rather than any architectural block. In a matched
nine-arm comparison, adding the OpenStereo augmentation triplet lowered the
best validation end-point error from 2.778 to 1.921 pixels, a 31 percent
reduction that is many times the run-to-run noise floor. The
edge-inference optimization pass, shown in Table
\ref{tab:ablation_efficiency}, held accuracy while cutting latency by
roughly a third: the exact-equivalence fixes and the narrow geometry
volume together reduced the latency from 46.6 to 30.2 milliseconds with
end-point error unchanged within noise, from 2.908 to 2.906 pixels. The
boundary-sharpness study of Table \ref{tab:ablation_blurfix} selected the
plane-rendering fix, which reduced the bad-0.5 rate from 55.36 to 42.91
percent, the sharpest result of the study. Importantly, the compositions
that added the bimodal head or the top-k initialization on top of plane
rendering were worse than plane rendering alone, at 44.40 and 46.58
percent bad-0.5, so the individual fixes anti-synergize and only plane
rendering was adopted. This negative result is reported because it
constrains future work. Finally, Table \ref{tab:ablation_native} shows the
input-protocol study on both evaluation axes: training on native crops cut
the native-axis end-point error from 6.672 to 2.869 pixels relative to the
legacy global-resize protocol, at no cost on the resized axis, which is
why the native-crop protocol was adopted for the production run.

\begin{table}[!htbp]
\centering
\caption{\textbf{Efficiency ablation.} 100-pair matched harness, 9k steps, identical seed and data; latency at $384\times 640$, batch 1.}
\label{tab:ablation_efficiency}
\small
\begin{tabular}{lccccc}
\toprule
Variant & Params (M) & EPE & bad-1 (\%) & D1-all (\%) & Latency (ms)\\
\midrule
base (gev4) & 2.962 & 2.908 & 30.32 & 16.31 & 46.6\\
+ exact-equivalence fixes & 2.962 & 3.021 & 32.09 & 17.79 & 35.8\\
+ narrow GEV band & 2.962 & 2.906 & 32.29 & 16.88 & 30.2\\
\bottomrule
\end{tabular}
\end{table}

\begin{table}[!htbp]
\centering
\caption{\textbf{Boundary-sharpness fixes and their composition.} 500-pair windowed split; the composition rows show that combining fixes is worse than plane rendering alone.}
\label{tab:ablation_blurfix}
\small
\begin{tabular}{lcccc}
\toprule
Variant & EPE & bad-0.5 (\%) & D1-all (\%) & Latency (ms)\\
\midrule
control (narrow) & 2.006 & 55.36 & 12.57 & 22.2\\
+ bundle1 (top-k init) & 1.982 & 49.95 & 12.14 & 21.9\\
+ bimodal aux head & 1.779 & 43.26 & 11.53 & 25.2\\
\textbf{+ plane rendering (adopted)} & 1.826 & \textbf{42.91} & 12.19 & 15.0\\
\midrule
plane + bimodal & 1.847 & 44.40 & 11.62 & 25.5\\
plane + bimodal + bundle1 & 1.853 & 46.58 & 11.93 & 17.2\\
\bottomrule
\end{tabular}
\end{table}

\begin{table}[!htbp]
\centering
\caption{\textbf{Input-protocol ablation.} Dual-axis final evaluation on the same 50 validation frames; each arm evaluated on both axes. Native crops win the native axis outright at no accuracy cost on the legacy resized axis.}
\label{tab:ablation_native}
\small
\begin{tabular}{lcccccc}
\toprule
& \multicolumn{3}{c}{Native axis ($960\times 540$)} & \multicolumn{3}{c}{Resized axis ($640\times 384$)}\\
\cmidrule(lr){2-4}\cmidrule(lr){5-7}
Arm & EPE & bad-1 & D1-all & EPE & bad-1 & D1-all\\
\midrule
global resize (legacy) & 6.672 & 58.69 & 35.81 & 1.783 & 29.91 & 11.44\\
\textbf{native crop (adopted)} & \textbf{2.869} & \textbf{36.75} & \textbf{16.47} & 1.783 & \textbf{29.30} & \textbf{10.71}\\
native full frame & 2.967 & 38.68 & 16.74 & 2.557 & 36.45 & 16.36\\
\bottomrule
\end{tabular}
\end{table}

\begin{figure}[!htbp]
\centering
\includegraphics[width=\textwidth]{fig_4_7_ablations}
\caption{\textbf{Ablation summary.} (a) OpenStereo augmentation triplet (matched 9-arm comparison, best validation EPE). (b) The edge-inference optimization pass (100-pair harness). (c) Boundary blur fixes and their compositions (500-pair split, bad-0.5). (d) Input protocol on the native evaluation axis.}
\label{fig:ablations}
\end{figure}

\section{Efficiency and Deployment Measurements}
Table \ref{tab:latency} reports the measured efficiency. On the RTX 3050
laptop GPU at batch one and 384 by 640 resolution, the optimized model ran
in 49.8 milliseconds in half precision and 61.4 milliseconds in full
precision, with an inference working set of 0.26 and 0.35 gigabytes
respectively. The target deployment device, a Jetson Orin Nano, was not
available for direct measurement, so the on-device figure is a projection
rather than a measurement. It is obtained from the measured half-precision
latency by applying three factors: an eager-to-TensorRT speed-up of 1.8
for a launch-bound network, a half-precision-to-eight-bit factor of 1.2,
and a device factor of 1.15 that averages the compute and bandwidth ratios
between the two GPUs. The projection places the model at roughly 30
milliseconds and 33 frames per second on the Orin Nano, at the real-time
target. These starred values will be replaced by direct measurements when
the hardware is available.

\begin{table}[!htbp]
\centering
\caption{\textbf{Latency and memory.} Measured values on RTX 3050 Laptop (batch 1, $384\times 640$, eager PyTorch); starred values are projections for Jetson Orin Nano INT8 TensorRT and are replaced by measurements when the device is available.}
\label{tab:latency}
\small
\begin{tabular}{lcccc}
\toprule
Configuration & Latency (ms) & FPS & Inference memory (GB)\\
\midrule
RTX 3050, fp32 & 61.4 & 16.3 & 0.35\\
RTX 3050, fp16 & 49.8 & 20.1 & 0.26\\
Jetson Orin Nano, INT8 TensorRT & ${\sim}30^{*}$ & ${\sim}33^{*}$ & ${\sim}0.15^{*}$\\
\bottomrule
\end{tabular}
\end{table}

\section{Qualitative Results}
The quantitative results are supported by visual inspection. On the
synthetic test set, Figure \ref{fig:sceneflow_qualitative} shows accurate
reconstruction of broad surfaces with error confined to occlusion
boundaries. On the unseen Middlebury scenes, Figure
\ref{fig:mb14_qualitative} shows that the easy scenes are reconstructed
cleanly and that failure is localized to thin repeated structures rather
than spread across the image. The plane-rendering step of Chapter 3 is
visible at boundaries, where the predicted disparity follows object edges
more closely than a bilinear upsampler would allow. Figure
\ref{fig:camera} extends the evidence to real data. Four indoor pairs from
the project's calibrated stereo rig, none of which the model was trained
or fine-tuned on, produce plausible disparity maps in which the foreground
subject is cleanly separated from the background, confirming that the
model transfers to a real camera without adaptation.

\begin{figure}[!htbp]
\centering
\includegraphics[width=\textwidth]{fig_4_6_camera}
\caption{\textbf{Real-camera qualitative results.} Four indoor pairs from the project's stereo rig ($1280\times 720$, rectified, no fine-tuning): left image and the model's predicted disparity at $384\times 640$. No ground truth exists for these pairs; the panels demonstrate zero-shot visual plausibility on real data.}
\label{fig:camera}
\end{figure}

\section{Societal, Environmental and Sustainability Aspects}
% REQUIRED by the OBE mapping (PO6, PO7).
The results carry practical implications beyond the benchmark numbers.
Running the depth computation on the device rather than in the cloud keeps
the video stream local, since no imagery need leave the sensor, which
matters for the privacy of the people and spaces a robot or headset
observes. It also removes the energy and latency of a network round trip:
a network operating within a power envelope of a few watts to a few tens
of watts consumes far less energy per frame than transmitting the video to
a datacenter and back. The choice of passive stereo over active sensing
has an environmental dimension as well. A stereo rig emits no laser or
structured light, uses two commodity image sensors rather than a
specialized time-of-flight or LiDAR unit, and therefore carries a lower
material cost and no emitted radiation. The training footprint is modest
and is reported transparently: the production model was trained in roughly
one day on a single cloud GPU, and the ablation and evaluation runs each
cost a fraction of a dollar of cloud time. A compact, low-cost, passive
depth sensor also lowers the barrier to robotics and assistive technology
in settings where specialized depth hardware is unaffordable.

\section{Discussion}
The results answer the four research questions of Chapter 3. The model
reaches 0.78 pixel end-point error in domain within a few million
parameters, which answers the first question in the affirmative. It
generalizes zero-shot to an unseen real domain at 10.9 percent D1-all,
far from the earlier chassis's collapse though short of the
distillation-trained references, which answers the second. It runs in
real time on the measured GPU and is projected to meet the target on the
intended device, which addresses the third pending on-device confirmation.
The fourth, tolerance to imperfect rectification, is addressed by the
perturbation study of Section \ref{sec:rectification}. Figure
\ref{fig:pareto_ours} places the model on the zero-shot
accuracy-efficiency front, between the earlier chassis and the larger
distillation-trained and iterative references. The cross-domain result
also speaks to an open question in the literature. The gap between the
earlier chassis and the present model was closed by the training pipeline,
not by adding iteration, and the strongest efficient reference in the
comparison, LiteAnyStereo, is itself non-iterative and owes its
cross-domain accuracy to a three-stage distillation recipe. The evidence
here is consistent with the view that the training pipeline, rather than
architectural iteration alone, carries cross-domain generalization,
although the present study does not isolate the two factors and the
question remains open. Several boundaries must be stated plainly. The
on-device Jetson latency is a projection, not a measurement. The results
come from a single training run, so run-to-run variance is not
characterized. The model has not been submitted to the KITTI leaderboard,
and it carries no explicit left-right consistency or occlusion reasoning.
These are the honest limits of the present work and they frame the future
directions of Chapter 5.

\begin{figure}[!htbp]
\centering
\includegraphics[width=0.78\textwidth]{fig_4_8_pareto_ours}
\caption{\textbf{StereoLite on the zero-shot accuracy-efficiency front.} Middlebury 2014 D1-all against trainable parameters; all four points share the identical 23-scene protocol of Table \ref{tab:mb14_comparison}.}
\label{fig:pareto_ours}
\end{figure}
